Many problems in statistics and machine learning can be formulated as the optimisation of a loss $\ell:\Theta\rightarrow\mathbb{R}$ whose gradient takes the form of an expectation with respect to an intractable parameter-dependent distribution $\pi_\theta$, i.e.\ $\nabla_\theta \ell(\theta) = \mathbb{E}_{\pi_\theta}[F(\theta, X)]$. This structure arises in problems such as maximum-likelihood training of \glspl*{ebm} \citep{cd_original,nijkamp2019learning,song2021train,oliva2025uniformintimeconvergenceboundspersistent}, latent \glspl*{ebm} \citep{pang2020learning,marks2025learninglatentenergybasedmodels}, \gls*{mmle} \citep{atchade2017perturbed,de_bortoli_efficient_2021,pmlr-v206-kuntz23a,lim2024momentumparticlemaximumlikelihood,akyildiz2025interacting}, fine tuning \citep{marion2025implicitdiffusionefficientoptimization}, and Bayesian experimental design \citep{iollo2025bayesian}.

These problems are usually solved via generic first-order optimisation schemes where each iteration involves running \gls*{mcmc} to sample from $\pi_\theta$, resulting in a nested optimisation-sampling scheme \citep{younes1999convergence,andrieu2014markovian,atchade2017perturbed,de_bortoli_efficient_2021}. However, this is computationally expensive with delicate convergence properties dependent on various factors such as Markov-chain bias, mixing, and step-size selections. A number of algorithms have been developed in recent literature to avoid separately equilibrating an inner Markov chain, see, e.g. \citet{pmlr-v206-kuntz23a,lim2024momentumparticlemaximumlikelihood,akyildiz2025interacting} for \gls*{mmle} in latent variable models, and \citet{marion2025implicitdiffusionefficientoptimization} for fine tuning of generative models. These methods construct joint sampling-optimisation procedures, typically with Langevin processes and gradient descent based optimisers. However, they are often tied to particular dynamical choices of optimisers and diffusions, which may even be replaced by alternatives in implementation. In fact, practical applications of such methods often use modern optimisers such as Adam \citep{kingma2017adammethodstochasticoptimization}, see, e.g., \citet{pmlr-v206-kuntz23a,marks2025learninglatentenergybasedmodels} or preconditioned Langevin dynamics, see, e.g., \citet{glyn2025statistical,wang2026training,marks2025learninglatentenergybasedmodels}. This raises the question of how to define principled, general purpose stochastic optimisation methods for loss functions with intractable gradients.

In this work, we develop general diffusion-based particle optimisation methods. Specifically:

\textbf{(C1)} We develop a general diffusion-based framework for stochastic optimisation of loss functions, accounting for general optimisation and sampling dynamics. This unifies and generalises several existing methods for general optimisation problems whose gradients involve parameter-dependent expectations and leads to the design of a novel class of algorithms.

\textbf{(C2)} Under this general formulation, we establish exponential contraction of the mean-field system under a joint strong-monotonicity condition and a finite-$N$ error bound for a particle approximation.

\textbf{(C3)} Using the general framework, we develop two novel algorithms, (i) a higher-order Langevin dynamics based \gls*{mmle} algorithm, and (ii) a momentum-based algorithm for training \glspl*{ebm}. We show that these algorithms outperform existing methods in their respective applications. We also discuss how to develop novel algorithms for other applications, by instantiating the general framework with appropriate choices of optimisers and diffusions.

The rest of the paper is organised as follows. In Section~\ref{sec:background}, we provide background on the optimisation problems we consider, and the sampling dynamics we use, together with a discussion of related work. In Section~\ref{sec:main_system}, we introduce our general mean-field system for stochastic optimisation. We then provide theoretical results on the convergence of the mean-field system and finite-particle error bounds in Section~\ref{sec:analysis}. Finally, we show the design of novel algorithms in Section~\ref{sec:algorithm} accompanied with numerical results and experiments.

\subsection{Notation}
We denote by $\mathcal{P}_2(\mathbb{R}^d)$ the set of probability measures on $\mathbb{R}^d$ with finite second moments.
%For a measure $\mu\in\mathcal{P}(\mathbb{R}^d)$ and a measurable function $f:\mathbb{R}^d\rightarrow\mathbb{R}$, we write $\mu(f) = \int f(x)\mu(\md x)$.
We write $\delta_x$ for the Dirac delta measure at $x$. For a function $f:\mathbb{R}^d\rightarrow\mathbb{R}$, we write $\nabla f$ for the gradient of $f$, and $\nabla^2 f$ for the Hessian of $f$. For a vector $v\in\mathbb{R}^d$, we write $\|v\|$ for the Euclidean norm of $v$. For a random vector $X$, $\Ltwonorm{X}:=(\mathbb{E}\|X\|^2)^{1/2}$. For a matrix $A\in\mathbb{R}^{d\times d}$, we write $\|A\|$ for the spectral norm of $A$. For two measures $\mu,\nu\in\mathcal{P}(\mathbb{R}^d)$, we write $W_2(\mu,\nu)$ for the 2-Wasserstein distance between $\mu$ and $\nu$:
\begin{align*}
    W_2(\mu,\nu) &= \left(\inf_{\gamma\in\Gamma(\mu,\nu)}\int \|x-y\|^2 \gamma(\md x,\md y)\right)^{1/2},
\end{align*}
where $\Gamma(\mu,\nu)$ is the set of couplings of $\mu$ and $\nu$. For a function $f:\mathbb{R}^d\rightarrow\mathbb{R}$, we write $\|f\|_\infty = \sup_{x\in\mathbb{R}^d}|f(x)|$ for the supremum norm of $f$. $L^p(\mu)$ denotes the space of functions where the integral $\int |f(x)|^p \mu(\md x)$ is finite.

We will denote various spaces belonging to optimiser, sampler, and auxiliary states using the letters $\sX, \sV, \Theta, \sM$. In all these cases, we assume that these spaces are Euclidean, i.e., $\sX = \mathbb{R}^{d_x}$, $\sV = \mathbb{R}^{d_v}$, $\Theta = \mathbb{R}^{d_\theta}$, $\sM = \mathbb{R}^{d_m}$. Therefore, their tangent spaces are also Euclidean spaces of the same dimensions. We will keep this setting for clarity and simplicity, and our results can be extended to more general spaces.